% sections/introduction.tex

\section{Introduction}
\label{sec:intro}

A self-modeling system is one that contains a faithful model of itself.
A companion paper~\cite{Paper5} argues that any finite-dimensional
system satisfying this definition necessarily has state spaces of the
form $M_n(\mathbb{C})^{\mathrm{sa}}$: self-modeling \emph{is} complex
quantum mechanics. The derivation proceeds through a chain of algebraic
forced moves: the sequential product structure of the model constrains
the state space to be a Euclidean Jordan algebra, and faithfulness of
the model then selects the complex matrix algebras.

A second companion paper~\cite{Paper7} argues that the exceptional Jordan
algebra $\hthree$ is the unique candidate for the ``basin'' in which
self-modeling occurs. The argument is non-composability: $\hthree$ is
the only simple Euclidean Jordan algebra that cannot be realized as a
tensor factor of a larger system~\cite{BarnumGraydonWilce2020}. An
observer (a C*-subalgebra of the form $M_n(\mathbb{C})^{\mathrm{sa}}$)
probing $\hthree$ via a rank-1 idempotent induces a Peirce decomposition
$27 = 1 \oplus 16 \oplus 10$. Paper~7 argues that the 16-dimensional
Peirce half-space $\Vhalf$ carries one generation of Standard Model
fermions, with chirality and the gauge group $\SMgauge$ emerging from
the observer's complex structure and the $\Ffour$ automorphism group.

The present paper is conditional on the results of Papers~5 and~7;
if those identifications hold, the results here follow.

This paper addresses the remaining sector: the 10-dimensional Peirce
complement $\Vzero = \htwo{O}$. We show that the algebraic data of
$\hthree$ uniquely match the exceptional entry in the
Gunaydin-Sierra-Townsend (GST) classification of $N{=}2$
Maxwell-Einstein supergravity theories~\cite{GST1983,GST1984}. The
identification proceeds in the same spirit as Paper~5: just as
the Jordan-von~Neumann-Wigner classification identifies the
self-modeling state spaces with complex quantum mechanics, the GST
classification identifies the algebraic data of $\hthree$ with the
exceptional $N{=}2$ MESGT.

\subsection{The problem}

Every known route to Einstein's equations from algebraic or information-theoretic
starting points requires substantial additional structure. Connes'
spectral action~\cite{ChamseddineConnes1997} assumes a manifold and a
Dirac operator. Jacobson's thermodynamic argument~\cite{Jacobson1995}
assumes local Lorentz invariance and a Rindler horizon structure.
Farnsworth's non-associative spectral geometry~\cite{Farnsworth2025}
works within $\hthree$ but has not been extended to include gravity.
All three approaches take a spacetime manifold as given.

Our advance is that the algebraic identity of spacetime is not assumed.
The Peirce complement $\Vzero$, projected through the C*-bottleneck to
$\htwo{C} \cong \mathbb{R}^{3,1}$, provides a 4-dimensional space with
Minkowski signature and the correct conformal algebra. We obtain the
\emph{dimension}, \emph{signature}, and \emph{symmetry algebra} of
spacetime from $\hthree$. The step from this algebraic Minkowski space
to a smooth manifold carrying dynamical fields is a shared assumption
with every quantum gravity program; we do not claim to derive it.
Given this assumption, the algebraic data of $\hthree$ satisfy all
three hypotheses of the GST classification theorem, and the GST
bijection uniquely identifies these data with the bosonic sector of
the exceptional $N{=}2$ MESGT, including the Einstein-Hilbert
term $-R/2$.

\subsection{Summary of results}

The identification chain is:

\begin{enumerate}[label=(\roman*)]
\item $\hthree$ is the self-modeling basin (Paper~7, non-composability).

\item An observer at a rank-1 idempotent $E$ induces the Peirce
  decomposition $\Vone \oplus \Vhalf \oplus \Vzero$ with dimensions
  $1 + 16 + 10$.

\item The C*-bottleneck projects $\Vzero = \htwo{O}$ to
  $\htwo{C}_u \cong \mathbb{R}^{3,1}$ with Minkowski signature
  $(+,-,-,-)$. The Kantor-Koecher-Tits construction gives
  $\KKT(\htwo{C}_u) = \mathfrak{so}(4,2)$, the 4-dimensional
  conformal algebra.

\item The unique $\Ffour$-invariant cubic form $\det(X)$ on $\hthree$
  (Springer, 1962) serves as the gravitational prepotential.
  Its role in $\rho_J$ and in gravity is forced by algebraic
  uniqueness (double duty corollary).

\item $\hthree$ satisfies all three GST hypotheses: degree~3, formally
  real, positive-definite trace form
  (Proposition~\ref{prop:gst-hypotheses}).

\item The GST bijection uniquely identifies the algebraic data of
  $\hthree$ with the exceptional 5-dimensional $N{=}2$ MESGT.

\item The $r$-map (dimensional reduction on a circle) gives the
  4-dimensional special K\"ahler theory with scalar manifold
  $\Eseven / (E_{6(-78)} \times \UU{1})$.

\item The $N{=}2$ label appears as the classification output: it is
  not assumed in any algebraic step, but enters through the GST
  theorem whose codomain is the class of $N{=}2$ MESGTs.
\end{enumerate}

The result is the complete bosonic Lagrangian
\begin{align}
\label{eq:lagrangian-preview}
e^{-1}\mathcal{L}_{\mathrm{bos}} &= -\frac{R}{2}
  + g_{i\bar{j}}\,\partial_\mu z^i \partial^\mu \bar{z}^{\bar{j}}
  \nonumber\\
  &\quad + \im(\mathcal{N}_{IJ})\, F^I_{\mu\nu} F^{J\mu\nu}
  + \re(\mathcal{N}_{IJ})\, F^I_{\mu\nu} {*F}^{J\mu\nu},
\end{align}
with 28 vectors (27~matter plus graviphoton), scalar manifold
$\Eseven / (E_{6(-78)} \times \UU{1})$, and all couplings determined
by $\det(X)$.

\subsection{Conventions}

We work in natural units ($\hbar = c = 8\pi G_N = 1$). The metric
signature is $(+,-,-,-)$. The Jordan product is
$\jp{X}{Y} = \tfrac{1}{2}(XY + YX)$. The octonion basis follows the
Fano plane convention $e_1 e_2 = e_4$, and the complex structure is
$u = e_7$ throughout. The Peirce idempotent is $E_{11}$
(the rank-1 projector onto the first diagonal entry of $\hthree$).
The real form of the structure group is $\Esix$ (not $E_{6(-78)}$ or
$E_{6(6)}$).
